\addcontentsline{toc}{section}{СПИСОК ИСПОЛЬЗОВАННЫХ ИСТОЧНИКОВ}

\begin{thebibliography}{99}
	
	%% ===== I. НОРМАТИВНЫЕ ДОКУМЕНТЫ =====
	
	\bibitem{gost19102} ГОСТ~19.102-77. Единая система программной документации. Стадии разработки. --- URL: \url{https://docs.cntd.ru/document/1200007652} (дата обращения: 07.05.2026).
	
	\bibitem{gost19105} ГОСТ~19.105-78. Единая система программной документации. Общие требования к программным документам. --- URL: \url{https://docs.cntd.ru/document/1200007655} (дата обращения: 07.05.2026).
	
	\bibitem{gost19201} ГОСТ~19.201-78. Единая система программной документации. Техническое задание. Требования к содержанию и оформлению. --- URL: \url{https://docs.cntd.ru/document/1200007659} (дата обращения: 07.05.2026).
	
	\bibitem{gost19301} ГОСТ~19.301-79. Единая система программной документации. Программа и методика испытаний. Требования к содержанию и оформлению. --- URL: \url{https://docs.cntd.ru/document/1200007663} (дата обращения: 07.05.2026).
	
	\bibitem{iso25010} ГОСТ~Р~ИСО/МЭК~25010-2015. Информационные технологии. Системная и программная инженерия. Требования и оценка качества систем и программного обеспечения (SQuaRE). Модели качества систем и программных продуктов. --- URL: \url{https://docs.cntd.ru/document/1200121069} (дата обращения: 07.05.2026).
	
	\bibitem{fz123} Федеральный закон от 22.07.2008 №~123-ФЗ «Технический регламент о требованиях пожарной безопасности». --- URL: \url{https://www.consultant.ru/document/cons_doc_LAW_78699/} (дата обращения: 07.05.2026).
	
	\bibitem{sp113130} СП~1.13130.2009. Системы противопожарной защиты. Эвакуационные пути и выходы. --- URL: \url{https://docs.cntd.ru/document/1200071143} (дата обращения: 07.05.2026).
	
	%% ===== II. ТЕХНИЧЕСКАЯ ДОКУМЕНТАЦИЯ =====
	
	\bibitem{json} Ecma International. The JSON Data Interchange Format. Standard ECMA-404, 2nd~ed. --- URL: \url{https://www.ecma-international.org/publications-and-standards/standards/ecma-404/} (дата обращения: 07.05.2026).
	
	\bibitem{maui} Microsoft. .NET Multi-platform App UI (MAUI) documentation. --- URL: \url{https://learn.microsoft.com/dotnet/maui/} (дата обращения: 07.05.2026).
	
	\bibitem{mauimvvm} Microsoft. Model-View-ViewModel (MVVM) in .NET MAUI. --- URL: \url{https://learn.microsoft.com/dotnet/maui/xaml/fundamentals/mvvm} (дата обращения: 07.05.2026).
	
	\bibitem{mauidi} Microsoft. Dependency injection in .NET MAUI. --- URL: \url{https://learn.microsoft.com/dotnet/maui/fundamentals/dependency-injection} (дата обращения: 07.05.2026).
	
	\bibitem{xaml} Microsoft. XAML overview in .NET MAUI. --- URL: \url{https://learn.microsoft.com/dotnet/maui/xaml/} (дата обращения: 07.05.2026).
	
	\bibitem{skiasharp} Microsoft. SkiaSharp graphics in .NET MAUI. --- URL: \url{https://learn.microsoft.com/dotnet/maui/user-interface/graphics/skiasharp/} (дата обращения: 07.05.2026).
	
	\bibitem{systemtextjson} Microsoft. System.Text.Json overview. --- URL: \url{https://learn.microsoft.com/dotnet/standard/serialization/system-text-json/overview} (дата обращения: 07.05.2026).
	
	\bibitem{navigineSdk} Navigine. Indoor Navigation and Positioning SDK Documentation. --- URL: \url{https://navigine.com/sdk-documentation/} (дата обращения: 07.05.2026).
	
	\bibitem{situmDocs} Situm Technologies. Situm Indoor Positioning Platform. Developer Documentation. --- URL: \url{https://situm.com/docs/} (дата обращения: 07.05.2026).
	
	\bibitem{mappedinDocs} Mappedin Inc. Mappedin Developer Documentation. --- URL: \url{https://developer.mappedin.com/docs/} (дата обращения: 07.05.2026).
	
	\bibitem{twogisFloors} 2ГИС. Поэтажные планы в MapGL JavaScript API. --- URL: \url{https://docs.2gis.com/ru/mapgl/reference/indoor} (дата обращения: 07.05.2026).
	
	\bibitem{mapstedOfficial} Mapsted Corp. Indoor Navigation Technology. --- URL: \url{https://mapsted.com/indoor-navigation/} (дата обращения: 07.05.2026).
	
	%% ===== III. РУССКОЯЗЫЧНЫЕ ИСТОЧНИКИ (39) =====
	
	\bibitem{amelichev2022} Амеличев~Г.Э., Панина~В.С. Сравнительный анализ алгоритмов поиска пути в помещении // E-Scio. --- 2022. --- URL: \url{https://cyberleninka.ru/article/n/sravnitelnyy-analiz-algoritmov-poiska-puti-v-pomeschenii} (дата обращения: 07.05.2026).
	
	\bibitem{cormen} Близнякова~Е.А., Куликов~А.А., Куликов~А.В. Сравнительный анализ методов поиска кратчайшего пути в графе // Архитектура, строительство, транспорт. --- 2022. --- Т.~1 (99). --- С.~80--87. --- DOI:~10.31660/2782-232X-2022-1-80-87.
	
	\bibitem{vinnikov2022} Виннников~М.Д. Разработка базы данных для мобильного приложения на операционной системе Android // Научные известия. --- 2022. --- URL: \url{https://cyberleninka.ru/article/n/razrabotka-bazy-dannyh-dlya-mobilnogo-prilozheniya-na-operatsionnoy-sisteme-android} (дата обращения: 07.05.2026).
	
	\bibitem{gof} Гамма~Э., Хелм~Р., Джонсон~Р., Влиссидес~Дж. Паттерны объектно-ориентированного проектирования / пер.~с~англ. --- СПб.~: Питер, 2021. --- 448~с. --- ISBN~978-5-4461-1595-2.
	
	\bibitem{gravit2019} Гравит~М.В., Карькин~И.Н., Дмитриев~И.И., Кузенков~К.А. Моделирование процесса эвакуации из высотных зданий и сооружений с использованием пассажирских лифтов // Пожаровзрывобезопасность. --- 2019. --- Т.~28, №~2. --- С.~66--80. --- DOI:~10.18322/PVB.2019.28.02.66-80.
	
	\bibitem{grinyak2018} Гриняк~В.М., Гриняк~Т.М., Цибанов~П.А. Позиционирование внутри помещений с помощью Bluetooth-устройств // Территория новых возможностей. Вестник ВГУЭС. --- 2018. --- Т.~10, №~2. --- С.~137--147. --- DOI:~10.24866/VVSU/2073-3984/2018-2/137-147.
	
	\bibitem{gusarova2023} Гусарова~А.А. Современные информационные технологии как инструмент повышения пожарной безопасности зданий // Пожаровзрывобезопасность. --- 2023. --- Т.~32, №~6. --- С.~36--46. --- DOI:~10.22227/0869-7493.2023.32.06.36-46.
	
	\bibitem{dijkstra} Гюллинг~А.О., Воронцова~Н.В. Визуализация алгоритма Дейкстры // Известия Тульского государственного университета. Технические науки. --- 2024. --- Т.~2. --- С.~127--128. --- DOI:~10.24412/2071-6168-2024-2-127-128.
	
	\bibitem{ershov2021} Ершов~Т.А. Обзор технологии мобильной разработки Xamarin // StudNet. --- 2021. --- URL: \url{https://cyberleninka.ru/article/n/obzor-tehnologii-mobilnoy-razrabotki-xamarin} (дата обращения: 07.05.2026).
	
	\bibitem{eskendir2019} Ескендир~М.А. Введение в разработку мобильных приложений // Вестник магистратуры. --- 2019. --- URL: \url{https://cyberleninka.ru/article/n/vvedenie-v-razrabotku-mobilnyh-prilozheniy} (дата обращения: 07.05.2026).
	
	\bibitem{zaitseva2021} Зайцева~И.А., Эбата~В.Ч., Ковбаса~Н.А. Практика применения методологий Agile, Scrum в ИТ-проектах // Индустриальная экономика. --- 2021. --- Т.~1. --- DOI:~10.47576/2712-7559\_2021\_1\_62.
	
	\bibitem{zhukovskaya2017} Жуковская~А.Н., Заушицина~А.С. Особенности разработки кроссплатформенных мобильных приложений // Решетнёвские чтения. --- 2017. --- URL: \url{https://cyberleninka.ru/article/n/osobennosti-razrabotki-krossplatformennyh-mobilnyh-prilozheniy} (дата обращения: 07.05.2026).
	
	\bibitem{levitin} Изотова~Т.Ю. Обзор алгоритмов поиска кратчайшего пути в графе // Новые информационные технологии в автоматизированных системах. --- 2016. --- URL: \url{https://cyberleninka.ru/article/n/obzor-algoritmov-poiska-kratchayshego-puti-v-grafe} (дата обращения: 07.05.2026).
	
	\bibitem{ilichev2018} Ильичёв~М.В., Мезенцева~Е.М. Фреймворк Xamarin // Теория и практика современной науки. --- 2018. --- URL: \url{https://cyberleninka.ru/article/n/freymvork-xamarin} (дата обращения: 07.05.2026).
	
	\bibitem{isaev2023} Исаев~М.И., Вахаева~Д.А., Алдамов~А.И. Алгоритмы для нахождения кратчайшего пути // Региональная и отраслевая экономика. --- 2023. --- №~4. --- С.~16--19. --- DOI:~10.47576/2949-1916\_2023\_4\_16.
	
	\bibitem{karmanova2020} Карманова~М.В. Концепция разработки мобильных приложений для картографирования зон чрезвычайных ситуаций и происшествий // Интерэкспо Гео-Сибирь. --- 2020. --- Т.~1, №~2. --- С.~68--74. --- DOI:~10.33764/2618-981X-2020-1-2-68-74.
	
	\bibitem{freeman} Кошелев~О.В. Шаблоны проектирования в программировании // Научный журнал. --- 2018. --- URL: \url{https://cyberleninka.ru/article/n/shablony-proektirovaniya-v-programmirovanii} (дата обращения: 07.05.2026).
	
	\bibitem{kulikov2021} Куликов~А.С. Обзор основных методов внутренней навигации // Форум молодых учёных. --- 2021. --- URL: \url{https://cyberleninka.ru/article/n/obzor-osnovnyh-metodov-vnutrenniy-navigatsii} (дата обращения: 07.05.2026).
	
	\bibitem{marinin2024} Маринин~А.К. Выбор архитектуры для мобильных приложений // Известия Кабардино-Балкарского НЦ РАН. --- 2024. --- Т.~26, №~5. --- С.~84--93. --- DOI:~10.35330/1991-6639-2024-26-5-84-93.
	
	\bibitem{fowler} Мартин~Р. Чистая архитектура. Искусство разработки программного обеспечения / пер.~с~англ. --- СПб.~: Питер, 2022. --- 352~с. --- ISBN~978-5-4461-0772-8.
	
	\bibitem{mongush2017} Монгуш~А.В., Кикин~П.М. Обзор технологий indoor-навигации // Интерэкспо Гео-Сибирь. --- 2017. --- URL: \url{https://cyberleninka.ru/article/n/obzor-tehnologiy-indoor-navigatsii} (дата обращения: 07.05.2026).
	
	\bibitem{naumov2021} Наумов~Р.К., Железков~Н.Э. Сравнительный анализ форматов хранения текстовых данных для дальнейшей обработки методами машинного обучения // Научный результат. Информационные технологии. --- 2021. --- Т.~6, №~1. --- С.~40--47. --- DOI:~10.18413/2518-1092-2021-6-1-0-5.
	
	\bibitem{ohotichenko2021} Охотиченко~А.В., Ухта~Ю.Б. Проектирование системы для навигации внутри здания со сложной иерархической структурой // Программные системы и вычислительные методы. --- 2021. --- №~4. --- С.~46--57. --- DOI:~10.7256/2454-0714.2021.4.37012.
	
	\bibitem{sazonov2024} Сазонов~А.П. Использование принципа Dependency Injection в мобильной разработке // Инновации и инвестиции. --- 2024. --- URL: \url{https://cyberleninka.ru/article/n/ispolzovanie-printsipa-dependency-injection-v-mobilnoy-razrabotke-na-territorii-rb} (дата обращения: 07.05.2026).
	
	\bibitem{samoshin2026} Самошин~Д.А., Георгиев~Н.А., Климовцова~Я.О., Кочетыгов~В.А. Затраты времени на изучение плана эвакуации и рекомендации по повышению его эффективности // Пожары и чрезвычайные ситуации: предотвращение, ликвидация. --- 2026. --- №~1. --- С.~26--36. --- DOI:~10.25257/FE.2026.1.26-36.
	
	\bibitem{albahari} Сергеев~К.С., Афонин~А.Ю., Осин~С.М. Особенности и различия архитектурных паттернов MV в языке программирования C\# // Парадигма. --- 2025. --- №~4. --- Ч.~2. --- URL: \url{https://cyberleninka.ru/article/n/osobennosti-i-razlichiya-arhitekturnyh-patternov-mv-v-yazyke-programmirovaniya-c} (дата обращения: 07.05.2026).
	
	\bibitem{sterlyagov2017} Стерлягов~С.П., Селютина~Е.П. Применение User Experience/User Interface моделирования для разработки мобильного приложения // Международный научно-исследовательский журнал. --- 2017. --- Т.~62. --- DOI:~10.23670/IRJ.2017.62.067.
	
	\bibitem{zhang2015} Ткачев~Д.В. Графы и их применение при чрезвычайных ситуациях // Вестник науки и образования. --- 2019. --- URL: \url{https://cyberleninka.ru/article/n/grafy-i-ih-primenenie-pri-chrezvychaynyh-situatsiyah} (дата обращения: 07.05.2026).
	
	\bibitem{richter} Троелсен~Э., Джепикс~Ф. Язык программирования C\#~7 и платформы .NET и .NET Core / пер.~с~англ. --- М.~: Диалектика, 2018. --- 1328~с. --- ISBN~978-5-6040723-1-8.
	
	\bibitem{uvarov2016} Уваров~А.Н. Инверсия управления и внедрение зависимостей // Символ науки. --- 2016. --- №~10-1. --- URL: \url{https://cyberleninka.ru/article/n/inversiya-upravleniya-i-vnedrenie-zavisimostey} (дата обращения: 07.05.2026).
	
	\bibitem{urbanovich2023} Урбанович~М.В., Ковалева~К.А. Принципы разработки пользовательских интерфейсов // Современные инновации, системы и технологии. --- 2023. --- Т.~3, №~4. --- DOI:~10.47813/2782-2818-2023-3-4-0363-0374.
	
	\bibitem{feshina2022} Фешина~Е.В., Пчела~Д.В., Долгов~А.В., Трофимов~Д.И. Анализ технологий разработки мобильных приложений и информационных систем на базе ОС Android // Вестник Адыгейского государственного университета. Сер.~4. --- 2022. --- Т.~1 (296). --- С.~85--91. --- DOI:~10.53598/2410-3225-2022-1-296-85-91.
	
	\bibitem{yarovaya2022} Яровая~Е.В. Принципы построения архитектуры программного обеспечения // E-Scio. --- 2022. --- URL: \url{https://cyberleninka.ru/article/n/printsipy-postroeniya-arhitektury-programmnogo-obespecheniya} (дата обращения: 07.05.2026).
	
	%% ===== IV. ИНОСТРАННЫЕ ИСТОЧНИКИ (11) =====
	
	\bibitem{deng2021} Deng~H., Ou~Z., Zhang~G., Deng~Y., Tian~M. BIM and Computer Vision-Based Framework for Fire Emergency Evacuation Considering Local Safety Performance // Sensors. --- 2021. --- Vol.~21, No.~11. --- Art.~3851. --- DOI:~10.3390/s21113851.
	
	\bibitem{han2020} Han~L., Guo~H., Zhang~H., Kong~Q., Zhang~A., Gong~C. An Efficient Staged Evacuation Planning Algorithm Applied to Multi-Exit Buildings // ISPRS International Journal of Geo-Information. --- 2020. --- Vol.~9, No.~1. --- Art.~46. --- DOI:~10.3390/ijgi9010046.
	
	\bibitem{jia2018} Jia~X., Ebone~A., Tan~Y. A Performance Evaluation of Cross-Platform Mobile Application Development Approaches // Proceedings of the 5th IEEE/ACM International Conference on Mobile Software Engineering and Systems (MOBILESoft~2018). --- Gothenburg, 2018. --- P.~92--93. --- DOI:~10.1145/3197231.3197252.
	
	\bibitem{landau2020} Landau~Y., Ben-Moshe~B. STEPS: An Indoor Navigation Framework for Mobile Devices // Sensors. --- 2020. --- Vol.~20, No.~14. --- Art.~3929. --- DOI:~10.3390/s20143929.
	
	\bibitem{mirahadi2020} Mirahadi~F., McCabe~B.Y. EvacuSafe: Building Evacuation Strategy Selection Using Route Risk Index // Journal of Computing in Civil Engineering (ASCE). --- 2020. --- Vol.~34, No.~2. --- Art.~04019051. --- DOI:~10.1061/(ASCE)CP.1943-5487.0000867.
	
	\bibitem{mirahadi2021} Mirahadi~F., McCabe~B.Y. EvacuSafe: A Real-Time Model for Building Evacuation Based on Dijkstra's Algorithm // Journal of Building Engineering. --- 2021. --- Vol.~34. --- Art.~101687. --- DOI:~10.1016/j.jobe.2020.101687.
	
	\bibitem{rubio2021} Rubio-Sandoval~J.I., Martínez-Rodríguez~J.-L., Lopez-Arevalo~I., Ríos-Alvarado~A.B., Rodríguez-Rodríguez~A.J., Vargas-Requena~D.T. An Indoor Navigation Methodology for Mobile Devices by Integrating Augmented Reality and Semantic Web // Sensors. --- 2021. --- Vol.~21, No.~16. --- Art.~5435. --- DOI:~10.3390/s21165435.
	
	\bibitem{singh2021} Singh~N., Choe~S., Punmiya~R. Machine Learning Based Indoor Localization Using Wi-Fi RSSI Fingerprints: An Overview // IEEE Access. --- 2021. --- Vol.~9. --- P.~127150--127174. --- DOI:~10.1109/ACCESS.2021.3111083.
	
	\bibitem{tomazic2021} Tomažič~S. Indoor Positioning and Navigation // Sensors. --- 2021. --- Vol.~21, No.~14. --- Art.~4793. --- DOI:~10.3390/s21144793.
	
	\bibitem{yen2021} Yen~H.-H., Lin~C.-H., Tsao~H.-W. Time-Aware and Temperature-Aware Fire Evacuation Path Algorithm in IoT-Enabled Multi-Story Multi-Exit Buildings // Sensors. --- 2021. --- Vol.~21, No.~1. --- Art.~111. --- DOI:~10.3390/s21010111.
	
	\bibitem{zhangfire2019} Zhang~J., Guo~J., Xiong~H., Liu~X., Zhang~D. A Framework for an Intelligent and Personalized Fire Evacuation Management System // Sensors. --- 2019. --- Vol.~19, No.~14. --- Art.~3128. --- DOI:~10.3390/s19143128.
	
\end{thebibliography}
